% !TeX program = xelatex
\documentclass[11pt,a4paper,UTF8]{ctexart}
\usepackage{geometry}
\geometry{left=2.5cm,right=2.5cm,top=2.5cm,bottom=2.5cm}
\usepackage{amsmath,amssymb,amsfonts,amsthm}
\usepackage{booktabs}
\usepackage{graphicx}
\usepackage{hyperref}
\usepackage{listings}
\usepackage{xcolor}
\usepackage{tcolorbox}
\usepackage{enumitem}

\hypersetup{
    colorlinks=true,
    linkcolor=blue!80!black,
    citecolor=blue!80!black,
    urlcolor=blue!80!black
}

\lstset{
    basicstyle=\small\ttfamily,
    keywordstyle=\color{blue}\bfseries,
    commentstyle=\color{gray}\itshape,
    stringstyle=\color{purple},
    frame=single,
    breaklines=true,
    numbers=left,
    numberstyle=\tiny\color{gray},
    xleftmargin=2em,
    framexleftmargin=1.5em
}

\title{\LARGE\textbf{数学建模平时作业}}
\author{\large 姓名：\underline{\hspace{3.5cm}} \qquad 学号：\underline{\hspace{3.5cm}}}
\date{\today}

\begin{document}
\maketitle

\begin{center}
\textbf{评分标准：问题重述(15\%) + 模型假设(15\%) + 建模与推导(35\%) + 求解与代码(25\%) + 分析评价(10\%)}
\end{center}
\vspace{0.5cm}

\section*{第 1 题}

\begin{tcolorbox}[title=\textbf{题目内容},colback=blue!5!white,colframe=blue!75!black,arc=2mm]
（在此填入教材或作业题目描述）
\end{tcolorbox}

\subsection*{一、 问题重述与机理分析}
针对实际背景，提炼核心物理量与决策目标。说明输入数据、因变量与自变量之间的基本依赖关系。

\subsection*{二、 模型假设与符号说明}
\begin{enumerate}[label=(\arabic*)]
    \item 假设系统处于稳态，忽略外界微弱扰动；
    \item 假设参数在考察周期内保持常数。
\end{enumerate}

\begin{table}[htbp]
\centering
\caption{核心参数与变量符号说明}
\vspace{1mm}
\begin{tabular}{cclc}
\toprule
符号 & 物理/经济含义 & 变量类型 & 单位 \\
\midrule
$t$ & 运行时间 & 自变量 & $\mathrm{s}$ \\
$x(t)$ & 系统核心状态量 & 状态变量 & $\mathrm{m}$ \\
$u(t)$ & 外部控制输入 & 决策变量 & --- \\
\bottomrule
\end{tabular}
\end{table}

\subsection*{三、 数学模型的建立与严密推导}
根据守恒定律或微积分极值原理，建立数学方程：
\begin{equation}
    \frac{\mathrm{d}x}{\mathrm{d}t} = f(x, u, t), \quad x(0) = x_0
\end{equation}

目标函数与边界约束条件为：
\begin{equation}
    \max_{u \in \mathcal{U}} \quad J = \int_0^T L(x(t), u(t)) \, \mathrm{d}t
\end{equation}

\subsection*{四、 模型求解与计算代码}
给出解析推导解或数值算法步骤。若附带 MATLAB/Python 求解代码：
\begin{lstlisting}[language=Python]
# 数学建模求解脚本示例
import numpy as np

def solve_model():
    pass
\end{lstlisting}

\subsection*{五、 计算结果分析、灵敏度分析与模型评价}
对求解得到的数值结果进行现实合理解释，讨论核心参数微扰对输出解的影响（灵敏度检验），指出模型的适用边界与改进方向。

\end{document}
